% Part: model-theory
% Chapter: interpolation
% Section: introduction

\documentclass[../../../include/open-logic-section]{subfiles}

\begin{document}

\olfileid{mod}{int}{int}

\olsection{مقدمة}

مبرهنة الاستيفاء هي النتيجة الآتية: افترض أن
$\Entails !A \lif !B$. عندئذ توجد !!a{sentence}~$!C$ بحيث
$\Entails !A \lif !C$ و$\Entails !C \lif !B$. وفوق ذلك، فإن كل
!!{constant} و!!{function} و!!{predicate} (عدا~$\eq$) في $!C$ يقع في
$!A$ و$!B$ معًا. وتسمى الـ!!{sentence}~$!C$ \emph{مستوفية} لـ$!A$
و$!B$.

لمبرهنة الاستيفاء أهمية في ذاتها، لكن أهميتها الرئيسية ترجع إلى إمكان
استعمالها لإثبات نتائج عن قابلية التعريف في نظرية، وعن الشروط التي يؤدي
فيها جمع نظريتين متسقتين إلى نظرية متسقة. تعرف النتيجة الأولى بمبرهنة بث
في قابلية التعريف، والثانية بمبرهنة روبنسون للاتساق المشترك.

\end{document}
